% =====================================================================
%  Section: The Particle Zoo   (Plan Stages 3-4)
%  Paper 1, "It from Bit via Gödel". Drafted 2026-06-08.
%  Follows \S sec:charges; precedes \S sec:masses.
%
%  \input fragment. biblatex/\autocite, amsthm `proposition'.
%  No new references (filatov2024towards, baez2002octonions,
%  gillardgresnigt2019three, furey2018three already in the library).
%
%  Cross-reference labels (placeholders -- re-point before compiling):
%    sec:charges               -- quantum numbers (previous section)
%    sec:entanglement-classes  -- GHZ/W theorems; Filatov chirality
%    sec:register              -- the three-qubit register; Filatov representation
%    prop:ghz, prop:w-doublet  -- the theorems
%    app:w-proof               -- the W proofs (Bell-shadow corollary)
%    sec:masses                -- masses (forward; mass hierarchy mechanism)
%    sec:interactions          -- circuits branch (forward; weak force as resolution-swap)
%
%  EPISTEMIC FLAGS carried through this section (per the source documents),
%  now recorded as named open problems (paper-wide convention, see the
%  quantum-numbers header for the preamble):
%    - op:filatov-extension  -- the three-qubit extension of the Filatov
%      constraint is ASSUMED, not proved: the foundational problem;
%    - op:chirality          -- the SU(2)_L-over-SU(2)_R fixed point;
%    - op:sterile-count      -- two-versus-three a-chiral modes;
%    - the mass ORDERING is structurally motivated, the mechanism deferred
%      to \S sec:masses (its open problems are converted when that section
%      is next worked).
%
%  REVISED 2026-06-11: open-problem convention implemented; four-class
%  table corrected (rows 2-3 had odd-qubit labels misaligned with the
%  handedness patterns); quark-doublet attribution corrected (frustrated-
%  pair resolutions, not the W-marginal spectrum, per the rescoped
%  prop:w-doublet); canonical-orbit scope; triality paragraph tied to the
%  explicit Fano lines of app:setup; NEW subsection "Why the colour force
%  is not chiral" (the Maudslay/Whitworth three-plates analogy --
%  historical attribution verified: Maudslay ~1800, Whitworth's 1840
%  British Association paper).
% =====================================================================

\section{The Particle Zoo}
\label{sec:generations}
% NOTE: revision of 2026-06-11 implements the open-problem convention and
% the review fixes; if a fuller rewrite follows, preserve: the label
% sec:zoo-singlets (anchor for the right-handed-singlet enumeration) and
% its two-way pairing with sec:rh-singlets in the quantum-numbers section;
% and the op: labels op:filatov-extension, op:chirality, op:sterile-count.

With a channel's quantum numbers now readable from its entanglement structure (\S\ref{sec:charges}),
we can enumerate the channels themselves: the fermion content of one generation, then why there are
exactly three of them, and finally a fourth class that couples to nothing. The spine of the argument
is a single topological fact --- the complete graph $K_3$ is not two-colourable --- applied to the
handedness that the Filatov representation forces on entangled qubits. It should be said at the
outset that one link in the chain is an assumption rather than a theorem, and a load-bearing one:
the extension of the Filatov constraint from two qubits to three. We flag it where it enters and
again at the close; the two-qubit result it generalises is proved \autocite{Filatov2024EntanglementStructure}, and
everything downstream is rigorous once the extension is granted.

\subsection{One generation}

A single generation is assembled from the two canonical entanglement orbits of
\S\ref{sec:entanglement-classes} --- the particle identification of \S\ref{sec:canonical-orbits} ---
together with chirality. The leptons are the W-class channels: a left-handed weak-isospin doublet,
the neutrino and the charged lepton, together with a right-handed charged-lepton singlet that the
weak force does not touch. The quarks are the GHZ-class channels, which carry colour; their
weak-isospin doublet does not come from the W-marginal spectrum --- a canonical GHZ orbit has
maximally mixed marginals and no distinguished axis --- but from the frustrated-pair resolutions
made precise below, the same $T_3 = \pm\tfrac12$ structure that Proposition~\ref{prop:w-doublet}
exhibits for the leptons. They appear as a left-handed doublet, the up- and down-type quark, in
three colours, with right-handed singlets. Collecting one generation:

\begin{center}
\begin{tabular}{lllc}
\hline
channel & class & multiplet & chirality \\
\hline
$(\nu,\ e)$            & W   & $SU(2)$ doublet, colour singlet      & left \\
$e$                    & W   & $SU(2)$ singlet, colour singlet      & right \\
$(u,\ d)\times 3$      & GHZ & $SU(2)$ doublet, colour triplet      & left \\
$u,\ d$ (each $\times3$) & GHZ & $SU(2)$ singlet, colour triplet    & right \\
\hline
\end{tabular}
\end{center}

The doublet is the two ways of resolving a single structure, made precise below; the singlet is the
channel with no such structure for the weak force to act on.\label{sec:zoo-singlets} The
right-handed singlets' hypercharges --- $Y = 2Q$ via the chirality flip --- are fixed in
\S\ref{sec:rh-singlets}, completing their quantum numbers. Nothing in the table is assigned by
hand --- the class fixes colour, the doublet fixes isospin, the chirality is the handedness the
next subsection forces, and the hypercharges are those of \S\ref{sec:charges}.

\subsection{The Filatov constraint and the \texorpdfstring{$K_3$}{K3} obstruction}

Chirality enters through the representation of \S\ref{sec:register}. Filatov and Auzinsh proved that
an entangled two-qubit state cannot be drawn on two Bloch spheres of the same handedness: one must be
right-handed and the other left-handed \autocite{Filatov2024EntanglementStructure}. This is the Filatov
constraint, and it is a topological consequence of the quaternionic Hopf fibration rather than a
drawing convention --- the handedness asymmetry is the geometric content of the entanglement itself.

A W-class channel has all three of its qubit pairs entangled: tracing any one qubit leaves the
genuine Bell pair of the bipartite shadow (Corollary in Appendix~\ref{app:w-proof}). Extending the
Filatov constraint to each pair independently, the three qubits must be assigned handednesses
$\chi_A,\chi_B,\chi_C \in \{L,R\}$ with every pair opposite,
\begin{equation}
  \chi_A \neq \chi_B, \qquad \chi_B \neq \chi_C, \qquad \chi_C \neq \chi_A.
\end{equation}
This is a request to two-colour the complete graph $K_3$ on three vertices. A graph is two-colourable
exactly when it has no odd cycle; $K_3$ \emph{is} an odd cycle, so it is not two-colourable. There
is no handedness assignment satisfying all three constraints at once. We call this the $K_3$
obstruction.

It must be stated plainly that the step just taken --- applying the constraint to each pair of a
\emph{tripartite} state independently --- is an assumption. The two-sphere theorem does not
formally imply it, and we record it as the foundational open problem of this section.

\begin{openproblem}[The Filatov extension]\label{op:filatov-extension}
Prove that each entangled pair of a tripartite state carries its own quaternionic Hopf structure,
so that the Filatov--Auzinsh opposite-handedness theorem \autocite{Filatov2024EntanglementStructure} applies to
every pair independently. The generation count of this section rests on this extension; the
combinatorics that follow are rigorous once it is granted.
\end{openproblem}

\subsection{Three generations and the doublet}

Since the three constraints cannot all hold, enumerate the $2^3 = 8$ handedness assignments and count
how many each satisfies. Every assignment satisfies either none or exactly two of the three (never
one, never three), giving four classes:

\begin{center}
\begin{tabular}{ccccll}
\hline
$\chi_A$ & $\chi_B$ & $\chi_C$ & satisfied pairs & odd qubit & reading \\
\hline
$\uparrow$ & $\downarrow$ & $\downarrow$ & $(AB),(AC)$ & $A$ ($\mathbb{C}$)  & generation 1 \\
$\uparrow$ & $\downarrow$ & $\uparrow$   & $(AB),(BC)$ & $B$ ($\mathbb{H}$)  & generation 2 \\
$\uparrow$ & $\uparrow$   & $\downarrow$ & $(AC),(BC)$ & $C$ ($\mathbb{O}$)  & generation 3 \\
$\uparrow$ & $\uparrow$   & $\uparrow$   & none        & ---                 & sterile (\S below) \\
\hline
\end{tabular}
\end{center}

Each row also has its globally conjugated partner ($\uparrow\!\leftrightarrow\!\downarrow$
throughout), which is the antiparticle. The first three classes each have exactly one qubit whose
handedness is the odd one out, and that qubit's Cayley--Dickson level labels the generation:
$\mathbb{C}$, $\mathbb{H}$, $\mathbb{O}$ for the first, second and third. The remaining
same-handedness pair --- the frustrated pair the constraint cannot satisfy --- has two topologically
distinct resolutions, and those two resolutions are the two members of the weak-isospin doublet,
$T_3 = \pm\tfrac12$. The weak force, in this reading, is the operation that switches between the two
resolutions of the frustrated pair, which is exactly the $W^\pm$ action on the $B$-qubit of
\S\ref{sec:charges}.

There is a deeper reason for the number three, beyond the combinatorial count, and it is worth
giving because it shows the count is not an artefact of the $K_3$ device. The three-qubit state
space carries the triality of $\mathrm{Spin}(8)$, which permutes the vector and the two spinor
representations $\mathbf{8}_v, \mathbf{8}_s, \mathbf{8}_c$ \autocite{baez2002octonions}. Fixing the
preferred complex direction breaks the triality and singles out the three Fano lines through that
direction --- the lines $(1,2,3)$, $(1,4,5)$ and $(1,7,6)$ of the multiplication table fixed in
Appendix~\ref{app:setup} --- each a quaternionic subalgebra; these are the three generations.
Triality gives three; the $K_3$ obstruction refines this to three-plus-one by adding the
unfrustrated class. The same count of three generations is reached independently from the complex
sedenions and from chains of division-algebra ideals \autocite{gillardgresnigt2019three,furey2018three},
a convergence we read as corroboration rather than as part of the present derivation.

The framework associates increasing mass with the increasing algebraic depth of the odd qubit ---
commutative $\mathbb{C}$ lightest, non-associative $\mathbb{O}$ heaviest --- which reproduces the
observed ordering $m_1 < m_2 < m_3$ in each charge sector. We flag honestly that what is settled here
is the \emph{count} and the depth \emph{ordering}; the mass values, and the mechanism that produces
the hierarchy, belong to \S\ref{sec:masses}, where it emerges that a naive associator count does not
by itself separate the generations and the hierarchy comes from a subtler feature of the geometry.
One prediction is immediate and sharp: there are four classes and only three are chiral, so there is
no fourth generation at any scale --- a stronger statement than the $Z$-width bound, and falsifiable.

\subsection{Why the colour force is not chiral}

A mechanical story from the dawn of precision engineering makes the structural point vivid. To
manufacture a truly flat surface, nineteenth-century engineers --- the practice originates with
Maudslay around 1800 and was formalised by Whitworth in 1840 --- would coat metal plates with
engineer's blue and rub them together: the blue reveals the high spots, which are scraped away, and
the cycle repeats. Done with \emph{two} plates the process converges, but not to flatness: it
produces a perfectly mating pair, one gently convex and the other gently concave --- a
complementary, handed pair whose error no two-plate comparison can ever detect, because convex fits
concave exactly. The cure was to work \emph{three} plates against one another pairwise, in
rotation. Writing each plate's deviation from flat as $\kappa$, the pairwise mating conditions
$\kappa_A = -\kappa_B$, $\kappa_B = -\kappa_C$, $\kappa_C = -\kappa_A$ force
$\kappa_A = \kappa_B = \kappa_C = 0$: the only surface that can mate with both of two others, which
also mate each other, is a plane. Two-way matching tolerates --- indeed manufactures --- a handed
pair; three-way cyclic matching destroys handedness.

This is the $K_3$ structure in continuous clothing, with one instructive difference. Handedness, in
the Filatov representation, is a property of \emph{pairs}: the two-sphere theorem assigns opposite
chirality to the members of an entangled pair, just as two-plate lapping assigns complementary
curvature. The weak force couples to exactly such a pair --- the frustrated pair, whose two
resolutions are the doublet --- so it is chirality-sensitive by construction; which chirality it
selects is Open Problem~\ref{op:chirality}, but \emph{that} it reads handedness follows from its
pairwise nature. The colour force couples instead to the three-way collective structure, the
cyclically matched triple, and a cyclically matched triple can carry no consistent handedness. There
is no handedness for a gate on the $C$-qubit to read, so colour couples to left and right alike ---
the one-generation table already exhibits this, with the right-handed singlets carrying colour
exactly as the left-handed doublets do --- and the same three-way symmetry leaves the QCD vacuum
angle with no pairwise orientation to inhabit (\S\ref{sec:entanglement-classes}). The instructive
difference is the resolution. Curvature is continuous, so the plates can satisfy all three
constraints at once by relaxing to the achiral fixed point $\kappa = 0$; binary handedness has no
such neutral value, so the qubit triple is \emph{frustrated} where the plates are \emph{flattened}
--- which is precisely why three chiral generations exist at all, each retaining one unsatisfiable
pair. The configuration that abandons pairwise opposition altogether, all three the same, is the
binary system's only achiral option, and it is the subject of the next subsection: the channel so
symmetric that no force has anything to grip.

\subsection{The fourth class: a sterile, a-chiral channel}

The fourth class has all three qubits of the same handedness. It satisfies none of the Filatov
constraints, so it has no opposite-handedness pair --- and with no such pair there is no $SU(2)$
structure for the weak force to act on. By the symmetry of the all-same configuration its net
hypercharge winding vanishes, so it is electrically neutral, and being a colour singlet it is inert
to the strong force as well. It couples to none of the gauge forces --- though whether it can truly
couple to \emph{nothing}, given that it gravitates, is a tension we return to below.

A terminological care is needed here, and it carries physical weight. This channel is not a
right-handed neutrino in the usual sense: that phrase presupposes chirality is a binary property
every fermion must carry, whereas the all-same configuration has \emph{no} opposite-handedness pair
and so is neither left- nor right-handed but \emph{a-chiral}. It is not the parity partner of the
active neutrino, which is why the standard objection to such a state as dark matter --- that parity
restoration at high energy should regenerate it --- does not apply: there is nothing for it to be the
partner of.

Two distinct notions of handedness must be kept apart at this point, because conflating them would
make ``a-chiral fermion'' sound like a contradiction. There is the \emph{external} chirality of the
Lorentz group --- whether a field sits in the $(\tfrac12,0)$ or $(0,\tfrac12)$ representation --- which
is bound up with spin, and there is the \emph{internal} handedness of the entanglement structure ---
the Filatov opposite-handedness on the $A,B,C$ register whose $K_3$ frustration generates the
generations. When we call this channel a-chiral we mean the \emph{internal} property: its handedness
assignment is the all-same $K_3$ class, the one carrying no opposite-handedness frustration for a
gauge force to engage. This says nothing against its spin. The channel is a three-qubit fermionic
register like every other fermion in the framework, and its spin-$\tfrac12$ character follows from
that, exactly as it does for the charged fermions; a-chirality is an additional internal property the
charged fermions possess and this one lacks, not the absence of the fermionic character that makes it
spin-$\tfrac12$ in the first place. Internal a-chirality is, moreover, precisely what makes it a
gauge singlet: with no opposite-handedness frustration there is nothing for $SU(2)$ to read, which is
the framework's derivation of the defining property of a sterile state --- that it does not feel the
weak force --- from the internal structure rather than as a label. Whether its Majorana mass is best
read as such, or admits a Dirac reading, we leave open.

The framework reads it as a sterile neutrino and a dark-matter candidate, with a
computability-scale Majorana mass driving a seesaw. On the count that two such a-chiral modes exist,
the seesaw yields a massless lightest active neutrino, $m_1 = 0$, normal mass ordering, and two
sterile states near a few TeV. The count of a-chiral modes is the weakest link in this part of the
chain, and we record it.

\begin{openproblem}[The sterile count]\label{op:sterile-count}
Determine, from the $K_3$ class structure, whether the a-chiral class supplies two or three
independent sterile modes to the seesaw. On the two-mode count the lightest active neutrino is
exactly massless, $m_1 = 0$, with normal ordering; on a three-mode count $m_1 = 0$ would not follow.
\end{openproblem}

A deeper tension attends this channel, and we state it at full strength rather than resolve it,
because the framework's own commitments make it sharp and its resolution would require a genuine new
insight. We said above that the channel couples to nothing; yet it carries a Majorana mass, and mass,
in this framework, is what gravitates (\S\ref{sec:confinement}). This is not the mild puzzle it would
be in non-relational physics, where a particle may simply interact through gravity alone. Here
gravity is not a force upon a pre-existing spacetime but the structure that measurement builds: the
framework identifies the act of measurement, the creation of a relation, the creation of spacetime,
and the gravitation we perceive as one process (\S\ref{sec:space-of-computations},
\S\ref{sec:fanout-boundary}). On that identification, whatever affects the causal structure must be
measurable --- must participate in relation-creating interactions --- and a channel that genuinely
couples to nothing cannot acquire a location or a causal role at all. A gauge-sterile massive state
is therefore in tension with the framework's own account: gravitational participation demands
measurable interaction, and gauge-sterility appears to deny it, while the usual escape, ``it
interacts only gravitationally,'' is unavailable precisely because gravity here is not a separate
interaction but the relational structure itself.

\begin{openproblem}[The sterile channel and relational gravity]\label{op:sterile-gravity}
Reconcile the gravitational participation of a gauge-sterile massive channel with the framework's
identification of gravity with measurement. Since whatever affects the causal structure must, on that
identification, be measurable, a channel that couples to no gauge force must nonetheless enter the
relational network by some route if it is to gravitate. Candidate directions, none here endorsed: that
the channel possesses a measurable but non-gauge interaction, which the framework would then have to
exhibit; that its mass --- its associator-geodesic length (Open Problem~\ref{op:associator-geodesic})
--- itself constitutes its relational participation, perhaps as a self-relation, which would connect
to the open Majorana reading and to the self-referential character of the framework as a whole; or
that the channel is a pre-spatial object that never acquires a distinct spacetime location, in which
case its cosmological gravitational effects must be accounted for without one. The resolution bears
directly on the nature of mass, of the sterile sector, and of dark matter, and we flag it as a
tension the framework's identification of gravity with measurement makes unavoidable.
\end{openproblem}

\subsection{What is settled, and what is owed}

Settled, resting on results proved earlier in the paper: colour from the canonical GHZ orbit
(Proposition~\ref{prop:ghz}); the lepton doublet from the canonical W orbit
(Proposition~\ref{prop:w-doublet}); the non-two-colourability of $K_3$, which is elementary; and,
granted the Filatov extension, exactly four classes, three chiral and one sterile, with the quark
doublet supplied by the frustrated-pair resolutions. The convergent three-generation counts from
triality and from the sedenions support the same conclusion.

Owed. One item has not yet been formalised: the selection of $SU(2)_L$ over $SU(2)_R$ --- why the
doublet is left-handed rather than right --- is a chirality fixed-point argument that propagates the
octonionic orientation down through the tower, sketched but not yet a theorem.

\begin{openproblem}[The chirality fixed point]\label{op:chirality}
Prove that the orientation of the octonionic level propagates down the Cayley--Dickson tower to
select $SU(2)_L$ over $SU(2)_R$ --- that the weak doublet is left-handed --- from the handedness
structure alone.
\end{openproblem}

The remaining debts are the problems already recorded: the Filatov extension
(Open Problem~\ref{op:filatov-extension}) is the foundation, and until it is proved the generation
count --- and with it the quark doublet --- is structurally motivated rather than derived; the
sterile count is Open Problem~\ref{op:sterile-count}. And the mass hierarchy, whose ordering the
algebraic depth motivates, awaits its mechanism in \S\ref{sec:masses}, where its open problems are
recorded in turn.

The spectrum is now enumerated: three generations of quarks and leptons in left-handed doublets and
right-handed singlets, with a fourth, sterile, a-chiral class. What the depth ordering asserts
qualitatively --- that mass grows with algebraic depth --- the next section makes quantitative,
turning the register geometry into the charged-lepton and quark mass spectra.
